% !TeX root = ../教材题库_全部条目_转写工作稿.tex
% Source: XJ-B1-C05-EX-P164-Y0865
% 第5章 / 5.2 任意角的三角函数 / 教材 p.164、165 / PDF 第171、172页
% OCR未校对：公式与文字均可能有误；未完成数学TeX转写；禁止直接发给学生
% 原题及相邻原解截图（核对用）：../题目原图/XJ-B1-C05-EX-P164-Y0865-01.webp, ../题目原图/XJ-B1-C05-EX-P164-Y0865-02.webp
% 本文件只在转写完成后作为题干使用；OCR 条目尚不是可靠题干。
\begin{bankitem}{5.2 例题 例题2}{第5章 / 5.2 任意角的三角函数 / 教材 p.164、165 / PDF 第171、172页}{XJ-B1-C05-EX-P164-Y0865}{OCR未校对：公式与文字均可能有误；未完成数学TeX转写；禁止直接发给学生}
\begin{ocrdraft}
[教材p.164 / PDF第171页]
\par
5丌7丌
\par
00分别求角
\par
一的正弦、余弦和正切值．
\par
4
\par
分析由于sina，cosa,tana的值与仅的终边上的点的位置无关，
\par
因而可以利
\par
用的终边与单位圆的交点来计算它们的值．
\par
解以单位圆的圆心0为原点建立平面直角坐标系，
\par
5丌
\par
召
\par
p
\par
5丌
\par
如图5．2一4．分别作Z，AOB一
\par
ZAOC\ocrglyph{00AF}—,设它们的
\par
6
\par
4
\par
终边与单位圆分别相交于点P()I，30，Q02，Y2).
\par
7丌
\par
4
\par
由于7一1，利用初中知识可得，点P,Q的坐标分别为
\par
2
\par
图5．2
\par
一4
\par
1
\par
因此，sin
\par
5丌
\par
5丌XI
\par
3
\par
5丌
\par
，COS
\par
6
\par
6
\par
7丌2
\par
同理可得sin
\par
2
\par
，COS
\par
4
\par
4
\par
4
\par
综上可知，对于每一个确定的角aeR,都分别有唯一确定的比值和与之
\par
对应，故正弦和余弦都是角a的函数．
\par
当角a的终边在3，轴上，也就是a一一十差丌（差ez）时，工=0，这时无意义．除
\par
[教材p.165 / PDF第172页]
\par
此之外，对于每一个确定的角a，都有唯一确定的比值与之对应，故正切也是角
\par
的函数．
\par
3，=仅，y=tana分别叫作角的正弦函数、余弦函数、正切函数．
\par
以上三种函数统称三角函数
\par
由于引人弧度制后，角的集合与实数集之间可以建立一一对应的关系，因此三
\par
角函数可以看成以实数为自变量的函数．在弧度制下，三角函数的定义域如下表：
\par
三角函数
\par
定义域
\par
此处的三鴦函数定义
\par
与初中所学的锐角三鴦函
\par
3，：=COS
\par
数有何异司？
\par
+差丌，差eZ
\par
30=tan
\end{ocrdraft}
\end{bankitem}
